
\begin{theorem}[Strong-CP closure on the admissible branch]
\label{thm:strong-cp-closure}
Define
\[
Z_{\mathrm{adm}}(\bar\theta)
:=
\int e^{-S_{\mathrm{adm}}}e^{i\bar\theta Q}\,d\mu_{\mathrm{adm}},
\qquad
Q=\frac{1}{32\pi^2}\int \tr(G\wedge G).
\]
On the admissible branch,
\[
\theta_{\mathrm{sheet}}=0,
\qquad
\arg\det M_u=\arg\det M_d=0,
\qquad
\bar\theta=0.
\]
Moreover,
\[
\bigl|Z_{\mathrm{adm}}(\bar\theta)\bigr|\le Z_{\mathrm{adm}}(0),
\]
so the free energy
\[
F(\bar\theta):=-V^{-1}\log Z_{\mathrm{adm}}(\bar\theta)
\]
satisfies $F(\bar\theta)\ge F(0)$. If the CP-even admissible vacuum is unique, its minimum lies at $\bar\theta=0$. At the same time weak CP may survive through
\[
V_{\mathrm{CKM}}=U_{u,L}^\dagger U_{d,L}.
\]
\end{theorem}

\begin{proof}[Proof sketch]
\Cref{thm:sheet-cp-protection,cor:strong-cp-rg-stability} supply the symmetry-protected vanishing of $\theta_{\mathrm{sheet}}$ and of the determinant phases on the retained branch. Positivity of the admissible measure then implies the standard free-energy inequality, so the CP-even branch minimizes at $\bar\theta=0$ once the admissible vacuum is unique. Weak CP may still survive through the relative family holonomies encoded in $V_{\mathrm{CKM}}$.
\end{proof}

\begin{theorem}[Reflection positivity on the admissible sector]
\label{thm:reflection-positivity-admissible}
Under the admissibility complex and the setup of \Cref{def:relative-generating-functional,lem:admissible-reflection-positivity}, every admissible observable $O$ with support in $\tau>0$ satisfies
\[
\langle \Theta O,O\rangle_{\mathrm{rel}}\ge 0.
\]
\end{theorem}

\begin{proof}[Proof sketch]
Invoke \Cref{lem:admissible-reflection-positivity}. The point of the quotient normalization is that the reference sector contributes only the constant denominator $Z_{\mathrm{ref}}[0,0,0]$ and therefore does not spoil positivity on admissible observables.
\end{proof}

\begin{theorem}[Lorentzian reconstruction on \texorpdfstring{$\Padm$}{Padm}]
\label{thm:lorentzian-reconstruction}
Under the same hypotheses, the admissible Euclidean theory reconstructs a self-adjoint Lorentzian Hamiltonian on the $\Padm$ sector.
\end{theorem}

\begin{proof}[Proof sketch]
With \Cref{thm:reflection-positivity-admissible} in hand, the bosonic part follows the usual Osterwalder--Schrader reconstruction, while the fermionic part follows the corresponding CAR reconstruction on the admissible sector. The commuting relations between $\Theta$, $P_\Theta$, and $\Padm$ ensure that the reconstructed Hilbert space remains inside the retained admissible branch.
\end{proof}

\begin{theorem}[Admissible vacuum existence at finite volume or with declared cutoff]
\label{thm:admissible-vacuum-existence}
Assume finite volume or a declared cutoff, weak-* compactness of $\Vadm$, and
lower boundedness plus coercivity of $\Padm\Hrel\Padm$. Then the admissible
vacuum exists as
\[
\rho^\ast=\argmin_{\rho\in\Vadm}\tr(\rho\Hrel),
\]
and its renormalized seam vacuum energy is denoted by $\Lambda_{\Seam,\mathrm{ren}}$.
\end{theorem}

\begin{proof}
The direct method applies on $\Vadm$ under the stated compactness and coercivity
assumptions. The renormalized seam residue supplies the vacuum-energy
identification in the same declared scheme.
\end{proof}

\begin{theorem}[Nonperturbative admissible QFT closure]
\label{thm:nonperturbative-admissible-qft}
Let
\[
m_*:=\min\left\{\frac{\epsilon_f}{2},\,m_-,\,m_3\right\}.
\]
Assume the transport gap estimate of \Cref{prop:gap-lift-full-operator} together with the positivity bounds of \Cref{thm:callias-higgs-selection}. Then, after removing the exact gauge zero modes, the reconstructed admissible Hamiltonian satisfies
\[
\operatorname{spec}(H_{\mathrm{adm}})\cap(0,m_*)=\varnothing.
\]
Hence connected admissible Euclidean correlators obey exponential clustering:
\[
\bigl|\langle O_1(x)O_2(0)\rangle_c\bigr|
\le
C_{12}e^{-m_*|x|}.
\]
\end{theorem}

\begin{proof}[Proof sketch]
The transport gap and the bosonic positivity bounds give a strict positive lower bound on all admissible excitations above the exact gauge zero modes. Reflection positivity and Lorentzian reconstruction convert the Euclidean gap into a Lorentzian spectral gap. Exponential clustering follows from that gap on the admissible sector.
\end{proof}

\begin{definition}[Stable record algebra from the admissible gap]
\label{def:stable-record-algebra}
Let
\[
\tau_*:=\frac{1}{m_*}.
\]
For a tolerance $\varepsilon>0$, define
\[
\Arec(t)
:=
\left\{
O\in\mathfrak A_{\mathrm{adm}}
:\ \|[O(t),O(t')]\|\le\varepsilon
\text{ for } |t-t'|>\tau_*
\right\}.
\]
\end{definition}

\begin{proposition}[Stable records from gap and clustering]
Under \Cref{thm:nonperturbative-admissible-qft}, the algebra $\Arec(t)$ defines the slow, robust sector singled out by the admissible mass gap. Record observables are therefore a dynamical consequence of gap plus clustering rather than an independent philosophical postulate.
\end{proposition}

\begin{proof}[Proof sketch]
The spectral gap suppresses long-time commutators by the same exponential scale that controls connected correlators. Therefore operators whose support is carried by the low-frequency admissible sector become approximately commuting on time separations larger than $\tau_*=1/m_*$, which is exactly the stability criterion built into \Cref{def:stable-record-algebra}.
\end{proof}

\begin{definition}[Observer algebra and prediction semantics]
\label{def:observer-algebra-prediction}
Choose $\Aobs\subset\Arec$ to be a maximal stable abelian subalgebra. For any operational query $Q$ represented by a spectral projector $E_Q(\Delta)$, define
\[
\Pred(Q):=\tr(\rho_tE_Q(\Delta)).
\]
This is the main-text prediction semantics used on the admissible branch; appendix-level bookkeeping only decorates it with extra experimental metadata.
\end{definition}

\begin{remark}[Horizons as modular inner seams]
The same seam grammar extends to horizon physics by treating horizons as modular inner seams. In that reading one works with algebraic splits or code-subspace factorizations adapted to the horizon sector, rather than declaring an exact tensor-product factorization as primitive structure.
\end{remark}

\section{Particle and mass compiler}

This section is the phenomenological heart of the manuscript. It is not meant to flatten all
rows into one uniform notion of prediction. Instead it makes the relation between derived theorem
outputs and intrinsic observation rows explicit: a declared numerical preconditioner, a
kernel-output map, and a mass ledger that separates internal theorem outputs from their
observation-surface images.
Accordingly this section is best read through three ledgers:
\begin{itemize}[leftmargin=1.5em, topsep=3pt, itemsep=2pt]
    \item a \emph{primary intrinsic-observation ledger} for the narrow main-text comparison surface in \Cref{tab:benchmark},
    \item \emph{mass-closure ledgers} for fermion, electroweak, and hadronic rows in \Cref{tab:sm-mass-ledger-ew,tab:sm-mass-ledger-quarks},
    \item and a \emph{supplementary appendix ledger} for exploratory or optional rows that are kept outside the closed theorem surface.
\end{itemize}

\begin{definition}[Intrinsic and observation surfaces]
\label{def:intrinsic-observation-surfaces}
We write $\mathcal S_{\mathrm{int}}$ for the intrinsic closure surface generated before
comparison, containing in particular the internal compiler masses $\hat m_f$, the transport
outputs, and the derived geometric parameters. The observation surface is
\[
\mathcal S_{\mathrm{obs}}:=\mathcal R_{\mathrm{TOE}}(\mathcal S_{\mathrm{int}}),
\]
that is, the scheme-specified comparison image after matching, thresholds, and observable
conventions are imposed. In practical finite-threshold computations one approximates this surface
through the numerical preconditioner $\mathcal R_{\SM}$. The first surface is theorem-facing; the
second is the observation-facing interface used by the ledgers below.
\end{definition}

\begin{figure}[p]
\centering
\scriptsize
\begin{tikzpicture}[
    scale=0.82,
    transform shape,
    >=Latex,
    node distance=9mm and 11mm,
    primitive/.style={
        draw=gray!70!black,
        thick,
        rounded corners,
        fill=gray!10,
        align=left,
        text width=3.55cm,
        inner sep=4pt
    },
    hard/.style={
        draw=blue!60!black,
        thick,
        rounded corners,
        fill=blue!5,
        align=left,
        text width=3.55cm,
        inner sep=4pt
    },
    bridge/.style={
        draw=green!50!black,
        dashed,
        thick,
        rounded corners,
        fill=green!5,
        align=left,
        text width=3.55cm,
        inner sep=4pt
    },
    closure/.style={
        draw=orange!75!black,
        dotted,
        thick,
        rounded corners,
        fill=orange!7,
        align=left,
        text width=3.65cm,
        inner sep=4pt
    },
    output/.style={
        draw=gray!70!black,
        thick,
        rounded corners,
        fill=gray!8,
        align=left,
        text width=3.45cm,
        inner sep=4pt
    },
    anchorrow/.style={
        draw=red!60!black,
        thick,
        rounded corners,
        fill=red!5,
        align=left,
        text width=3.45cm,
        inner sep=4pt
    },
    flow/.style={->, thick},
    aux/.style={->, thick, dashed},
    cond/.style={->, thick, dotted}
]

\node[primitive] (core) {%
\textbf{Primitive core}\\
$(X,S,\Hint,D_{\mathrm{rel}},\iota,f,\chiseed)$\\
undecomposed spectral seed
};

\node[hard, below=of core] (carrier) {%
\textbf{Minimal carrier kernel}\\
$E_3 \oplus E_2$\\
$Y=(-1/3)^3\oplus(1/2)^2$\\
$6Y^2-Y-\mathbf 1=0$
};

\node[hard, below left=16mm and 13mm of carrier] (familypacket) {%
\textbf{Hard packet}\\
$G_{\mathrm{SM}}^{\mathrm{glob}} \cong S(U(3)\times U(2))$\\
$S^+=\Lambda^{\mathrm{even}}E$\\
$S^+ \cong \mathbf 1 \oplus \mathbf{10} \oplus \overline{\mathbf 5}$
};

\node[hard, below right=16mm and 13mm of carrier] (seed) {%
\textbf{Hard kernel formulas}\\
$c_3=\dfrac{1}{8\pi}$,\quad
$\delta_{\mathrm{top}}=\dfrac{3}{256\pi^4}$\\
$\phi_{\mathrm{EM}}(\alpha)=\phibase+\deltatop x(1-\deltatop x)^{-5/4}$,\ $x=e^{-2\alpha}$\\
$b_1(N_{\mathrm{fam}},N_\Phi)=\dfrac43N_{\mathrm{fam}}+\dfrac1{10}N_\Phi$
};

\node[closure, below=of familypacket] (familybridge) {%
\textbf{Family closure}\\
$\Padm$, $F\cong\mathbb C^3$\\
$\mathrm{Ind}_{APS}(D^+_{\mathrm{rel}}\!\otimes\!L_F)=3$\\
$N_{\mathrm{fam}}=3$, $\Omega_{\mathrm{adm}}=48$
};

\node[closure, below=of seed] (chi) {%
\textbf{Local spectral-scale / Higgs branch}\\
$\mathrm{Ind}^{\Padm}_{\mathrm{rel}}(\mathcal B^+_{E_2})=1$\\
$\bar M_{\mathrm{Pl}}^2=\chigeozero^2/(2\pi^2)$\\
$\vgeo^2=Z_H\mu_\Phi^2(\chigeozero)/\lambda_\Phi$
};

\node[bridge, right=20mm of familybridge] (comparison) {%
\textbf{Numerical preconditioner $\mathcal R_{\mathrm{SM}}$}\\
matching scale, thresholds, scheme map\\
first observation iterate for $\mathcal R_{\mathrm{TOE}}$
};

\node[closure, below=of familybridge] (transport) {%
\textbf{Transport compiler}\\
$Y_f=D_{L,f}U_fD_{R,f}$\\
$\lambda_Y=\sqrt{\phiz(1-\phiz)}$, $D_y$, $U_6$\\
$(Y_f)_{ij}\sim \lambda_Y^{L^L_{f,i}+L^R_{f,j}}\mathrm{Hol}_{\Gamma_{ij}}$
};

\node[output, below right=12mm and 12mm of transport] (massmatrix) {%
\textbf{Observation lift}\\
$M_f=\dfrac{\vgeo}{\sqrt 2}\,Y_f$\\
$Y_f=D_{L,f}\,U_f\,D_{R,f}$\\
$V_{\mathrm{CKM}}=U_{u,L}^\dagger U_{d,L}$,\quad
$U_{\mathrm{PMNS}}=U_{\nu,L}^\dagger U_{e,L}$
};

\node[bridge, right=24mm of chi] (benchmarks) {%
\textbf{Intrinsic observation rows}\\
$\alpha^{-1}(0)$,\ $\beta_{\mathrm{rad}}$,\ $\Omega_b$\\
$\lambda_C$,\ $\sin^2\theta_{13}$\\
comparison-facing rows under $\mathcal R_{\mathrm{TOE}}$
};

\node[anchorrow, below=of benchmarks] (bosons) {%
\textbf{Gauge / Higgs masses}\\
exact: $m_\gamma=0$, $m_g=0$\\
intrinsic observation: $m_W$, $m_Z$, $m_H$\\
global threshold outputs
};

\node[output, below left=11mm and 12mm of massmatrix] (leptons) {%
\textbf{Charged leptons}\\
derived rows:\\
$m_e,\ m_\mu,\ m_\tau$
};

\node[anchorrow, below=of massmatrix] (quarks) {%
\textbf{Quark sector}\\
derived: $m_u,\ m_d$\\
intrinsic observation rows: $m_s,\ m_c,\ m_b,\ m_t$
};

\node[output, below right=11mm and 12mm of massmatrix] (neutrinos) {%
\textbf{Neutrino block}\\
$m_{\nu_1},\ m_{\nu_2},\ m_{\nu_3}$\\
closed neutrino branch + observation map\\
$M_R \sim 1.31\times 10^{15}\,\mathrm{GeV}$
};

\node[output, below=of quarks] (hadrons) {%
\textbf{Hadronic checks}\\
confinement as admissible singlet selection\\
$m_{\pi^0}$,\ $m_p$ as out-of-sample checks
};

\node[closure, below=of bosons] (cp) {%
\textbf{Closure outputs}\\
$\bar\theta=0$,\quad $\langle\Theta O,O\rangle_{\mathrm{rel}}\ge 0$\\
Lorentzian reconstruction\\
vacuum minimizer $\rho^\ast$, $\Lambda_{\Seam,\mathrm{ren}}$
};

\draw[flow] (core) -- (carrier);
\draw[flow] (carrier) -- (familypacket);
\draw[flow] (carrier) -- (seed);
\draw[flow] (familypacket) -- (familybridge);
\draw[aux] (seed) -- (familybridge);
\draw[flow] (seed) -- (comparison);
\draw[flow] (seed) -- (benchmarks);
\draw[flow] (seed) -- (transport);
\draw[aux] (familybridge) -- (transport);
\draw[flow] (chi) -- (bosons);
\draw[flow] (chi) -- (massmatrix);
\draw[flow] (transport) -- (massmatrix);
\draw[aux] (comparison) -- (benchmarks);
\draw[flow] (massmatrix) -- (leptons);
\draw[flow] (massmatrix) -- (quarks);
\draw[flow] (massmatrix) -- (neutrinos);
\draw[aux] (transport) -- (hadrons);
\draw[aux] (quarks) -- (hadrons);
\draw[cond] (transport) -- (cp);
\draw[cond] (chi) -- (cp);

\begin{scope}[on background layer]
    \node[draw=gray!55!black, rounded corners, inner sep=7pt,
          fit=(core)] {};
    \node[draw=blue!30!black, rounded corners, inner sep=7pt,
          fit=(carrier)(familypacket)(seed)] {};
    \node[draw=green!35!black, dashed, rounded corners, inner sep=7pt,
          fit=(comparison)(benchmarks)] {};
    \node[draw=orange!55!black, dotted, rounded corners, inner sep=7pt,
          fit=(familybridge)(chi)(transport)(cp)] {};
    \node[draw=gray!55!black, rounded corners, inner sep=7pt,
          fit=(bosons)(massmatrix)(leptons)(quarks)(neutrinos)(hadrons)] {};
\end{scope}
\end{tikzpicture}
\caption{Carrier-first compiler from reduced datum to explicit particle-content and mass
surfaces. Gray nodes denote the reduced datum, blue nodes are exact algebraic carrier/kernel
statements, orange nodes are conditional closure statements, green nodes denote the matched
intrinsic observation layer and its numerical preconditioner, and gray/red nodes are compact
output surfaces. The figure intentionally separates exact algebraic and conditional theorem
outputs from matched observation rows.}
\label{fig:compiler-to-masses}
\end{figure}

The diagram in \Cref{fig:compiler-to-masses} is epistemically layered: it visualizes one
compiler flow, not one uniform status class. Exact massless gauge directions, conditional
charged-lepton rows, matched heavy-mass observation rows, neutrino observation rows, and
conditional hadronic rows are shown together only to display the architecture, not to erase
their different proof roles.

A second, more compressed projection of that same surface is shown in \Cref{fig:epistemic-output-matrix}. It does not add a new theorem; it makes the manuscript's status discipline readable at the level of output types.
